\section{Introduction}

Verified software does not remain static. As code evolves, specifications, assertions, invariants, and proof hints must evolve with it. In mainstream software engineering, refactoring is a routine part of maintenance rather than an exceptional event~\cite{Mens2004RefactoringSurvey}, and empirical studies show that real projects perform a wide range of such structural changes during ordinary development~\cite{Kim2012RefactoringInTheWild,Tsantalis2014RefactoringOperations}. More broadly, maintenance research has emphasized that software artifacts are best understood through co-evolution across versions rather than as isolated one-shot tasks~\cite{Zaidman2010CoEvolutionProductionTest}. For verified programs, however, one basic empirical question remains underexplored: how brittle is verification under semantics-preserving refactoring?

This question is especially salient in auto-active verification settings such as Dafny, where program text and proof-relevant structure are intertwined in a single workflow~\cite{Leino2010Dafny}. Dafny automates much of proof checking, but successful verification still depends on carefully arranged contracts, loop invariants, assertions, and ghost code~\cite{Grov2016VerificationPatternsDafny}. Recent studies of Dafny practice suggest that maintainability and workflow friction remain central to the practical cost of verification adoption~\cite{Mugnier2025ImpactVerification}. The issue, then, is not proving from scratch, but understanding how nearby-version edits disrupt proofs that previously discharged automatically.

Existing literature only partially addresses this problem. Refactoring is well studied in software engineering, and proof refactoring has been explored in interactive theorem proving~\cite{Whiteside2011ProofScriptRefactoring,Whiteside2013RefactoringProofs,Adams2015Tactician}. The Dafny ecosystem has also seen recent work on benchmarks, testing, and learning-based support~\cite{DafnyBench2024,Wang2026TACoDafny,Xu2026DafnyComp}. What remains missing is a reproducible benchmark and empirical characterization of \emph{semantics-preserving refactoring-induced proof breakage in already verified Dafny programs}. This paper targets that gap.

To do so, it introduces Dataset~A, a maintenance-centred benchmark of verified Dafny programs paired with transformed variants drawn from seven refactoring classes. The benchmark is used to measure how often verification breaks, how failures cluster by proof-relevant failure mode, and how much breakage simple baseline recovery can repair. The baselines are informed by automated repair and proof-automation work~\cite{Gazzola2017AutomaticSoftwareRepair,Xia2023APRWithPLM,First2020TacTok,SanchezStern2023Passport,Misu2024AIAssistedDafny}, but the paper is not itself a repair-method contribution. Its central claim is narrower and, we argue, foundational: nearby-version proof breakage is a measurable maintenance problem, and lightweight recovery strategies leave a substantial residual repair gap.

Our contributions are:
\begin{itemize}
  \item a reproducible Dafny benchmark for refactoring-induced proof brittleness in already verified programs;
  \item an empirical characterization of breakage rates and failure-mode clusters across seven semantics-preserving refactoring classes;
  \item a baseline comparison of re-verification, syntactic transfer, and LLM-only repair that quantifies the residual repair gap while also providing an evaluation substrate for later repair work.
\end{itemize}

The rest of the paper develops this argument in a compact sequence. Section~2 clarifies the conceptual framing needed to separate nearby-version proof brittleness from other forms of change, Section~3 describes Dataset~A and the refactoring benchmark, Section~4 specifies the baselines and evaluation setup, and Section~5 presents the empirical results and their intended interpretation. Section~6 discusses threats to validity, Section~7 positions the paper against related work, and Section~8 concludes by summarizing why the benchmark is a standalone contribution as well as a foundation for later repair work.
